% ============================================================
%  preamble.tex — styling, environments, macros
%  Edit this file to change the look of the whole book.
% ============================================================

% ---------- Page geometry ----------
\usepackage[
  paperwidth=17cm, paperheight=24cm,
  inner=2.2cm, outer=1.8cm, top=2.2cm, bottom=2.4cm,
  headsep=0.7cm, footskip=1.2cm
]{geometry}

% ---------- Encoding & fonts ----------
\usepackage[T1]{fontenc}
\usepackage[utf8]{inputenc}
\IfFileExists{lmodern.sty}{\usepackage{lmodern}}{}
% Nicer fonts if available on your machine (MiKTeX / full TeX Live).
% Harmless no-ops if the packages are missing.
\IfFileExists{newtxtext.sty}{\usepackage{newtxtext}}{}
\IfFileExists{newtxmath.sty}{\usepackage{newtxmath}}{}
\IfFileExists{inconsolata.sty}{\usepackage[varqu,scaled=0.95]{inconsolata}}{}
% Font expansion needs scalable fonts; fall back gracefully if absent.
\IfFileExists{lmodern.sty}{\usepackage{microtype}}{\usepackage[expansion=false]{microtype}}

% Appendix E is a bilingual glossary. Polish needs T1 (set above) plus its
% hyphenation patterns, which live in a separate TeX Live language collection
% that a minimal installation does not have. Loading babel with a language whose
% .ldf is absent is a *fatal* error, not a warning, so both the package and the
% language file are probed before either is requested. British is likewise
% optional: without it the book still sets, it just hyphenates as American.
\IfFileExists{babel.sty}{%
  \IfFileExists{polish.ldf}{\usepackage[polish,main=british]{babel}}{%
    \IfFileExists{british.ldf}{\usepackage[british]{babel}}{}%
  }%
}{}

% ---------- Core packages ----------
\usepackage{graphicx}
\usepackage{xcolor}
\usepackage{booktabs}
\usepackage{tabularx}
\usepackage{longtable}
\usepackage{enumitem}
\usepackage{listings}
\usepackage[most]{tcolorbox}
\usepackage{fancyhdr}
\usepackage{titlesec}
\usepackage{caption}
\usepackage{etoolbox}
\usepackage{imakeidx}
\usepackage[hidelinks]{hyperref}

\hypersetup{
  colorlinks=true,
  linkcolor=lcblue,
  urlcolor=lcblue,
  citecolor=lcblue,
  pdftitle={LangChain, LangGraph and Async Python},
  pdfauthor={Konrad Cinkusz}
}

% ---------- Palette ----------
\definecolor{lcblue}{HTML}{1F4E79}
\definecolor{lcteal}{HTML}{0E7C7B}
\definecolor{lcamber}{HTML}{B26A00}
\definecolor{lcred}{HTML}{9B2C2C}
\definecolor{lcgrey}{HTML}{5A5A5A}
\definecolor{lclight}{HTML}{F4F6F8}
\definecolor{codebg}{HTML}{FAFAFA}
\definecolor{codeframe}{HTML}{D8DEE4}
\definecolor{kw}{HTML}{0B5394}
\definecolor{str}{HTML}{9B2C2C}
\definecolor{cmt}{HTML}{4F7A4F}
\definecolor{typ}{HTML}{0E7C7B}
\definecolor{dec}{HTML}{7B4397}

% ---------- Headings ----------
% \raggedright on the title body is load-bearing, not cosmetic. Several chapter
% titles here run past 35 characters, and at \Huge on a 13 cm text block a title
% that cannot break overflows the margin by 30 pt or more. Justified setting
% gives TeX no way out; ragged right lets it wrap onto a second line.
\titleformat{\chapter}[display]
  {\normalfont\huge\bfseries\color{lcblue}\raggedright}
  {\normalfont\normalsize\color{lcgrey}\MakeUppercase{\chaptertitlename}\ \thechapter}
  {8pt}{\Huge\raggedright}
\titlespacing*{\chapter}{0pt}{10pt}{28pt}
\titleformat{\section}{\normalfont\Large\bfseries\color{lcblue}}{\thesection}{0.8em}{}
\titleformat{\subsection}{\normalfont\large\bfseries\color{lcgrey}}{\thesubsection}{0.7em}{}

\pagestyle{fancy}
\fancyhf{}
\fancyhead[LE]{\small\nouppercase{\leftmark}}
\fancyhead[RO]{\small\nouppercase{\rightmark}}
\fancyfoot[LE,RO]{\small\thepage}
\renewcommand{\headrulewidth}{0.4pt}

% The default book-class running head is "Chapter 10. <full title>", and at this
% page width the longer chapter titles in Parts III and IV overflow the margin by
% 20-34 pt. Dropping the word "Chapter" recovers roughly that much and costs the
% reader nothing: the number is still there, and the head stays unambiguous.
%
% This MUST come after \pagestyle{fancy}: fancyhdr installs its own \chaptermark
% at that point and silently overwrites anything defined earlier.
\renewcommand{\chaptermark}[1]{\markboth{\thechapter.\ #1}{}}
\renewcommand{\sectionmark}[1]{\markright{\thesection\ #1}}

% ============================================================
%  CODE LISTINGS
% ============================================================

% `listings` ships a Python definition, but it knows nothing about the
% vocabulary this book actually uses. Extending it means `await`, `async`,
% `match` and the library's own type names are highlighted rather than left
% as plain identifiers in the middle of a listing.
\lstdefinelanguage{pythonx}{
  language=Python,
  morekeywords={
    async,await,match,case,nonlocal,yield,
    TypedDict,Annotated,Protocol,Literal,Generic,Final,ClassVar,
    BaseModel,Field,
    StateGraph,START,END,Command,Send,Runtime,
    RetryPolicy,CachePolicy,Interrupt,
    AIMessage,HumanMessage,SystemMessage,ToolMessage
  },
  morekeywords=[2]{
    create_agent,init_chat_model,interrupt,add_messages,tool,entrypoint,task,
    ainvoke,astream,astream_events,abatch,
    gather,create_task,TaskGroup,to_thread,timeout,wait_for,Semaphore,Queue
  },
  sensitive=true
}

\lstdefinestyle{lccode}{
  basicstyle=\ttfamily\footnotesize,
  keywordstyle=\color{kw}\bfseries,
  keywordstyle=[2]\color{typ},
  stringstyle=\color{str},
  commentstyle=\color{cmt}\itshape,
  emphstyle=\color{dec}\bfseries,
  numbers=left,
  numberstyle=\ttfamily\tiny\color{lcgrey},
  numbersep=8pt,
  backgroundcolor=\color{codebg},
  frame=single,
  rulecolor=\color{codeframe},
  framesep=6pt,
  breaklines=true,
  breakatwhitespace=false,
  postbreak=\mbox{\textcolor{lcgrey}{$\hookrightarrow$}\space},
  showstringspaces=false,
  tabsize=4,
  columns=fullflexible,
  keepspaces=true,
  captionpos=b,
  aboveskip=1.2em,
  belowskip=1.2em,
  % Maps the characters most likely to be pasted into a listing by accident.
  %
  % Do NOT add a mapping for U+00A0 (non-breaking space) here. It looks like the
  % obvious next entry and it makes `listings` abort the run with "Improper
  % alphabetic constant" -- a fatal error that names neither the character nor
  % this line. The ASCII-in-listings convention is what guards against it
  % instead; see CLAUDE.md.
  literate={ł}{{\l}}1 {Ł}{{\L}}1 {ą}{{\k a}}1 {ę}{{\k e}}1
           {ś}{{\'s}}1 {ć}{{\'c}}1 {ż}{{\.z}}1 {ź}{{\'z}}1 {ó}{{\'o}}1 {ń}{{\'n}}1
           {—}{{-{}-}}2 {–}{{-}}1 {‘}{{'}}1 {’}{{'}}1
           {“}{{"}}1 {”}{{"}}1 {°}{{\textdegree}}1
}

\lstset{style=lccode}

% Inline Python environment.  Usage:
%   \begin{python}[caption={Hello agent},label={lst:hello}]
%   ...code...
%   \end{python}
\lstnewenvironment{python}[1][]
  {\lstset{language=pythonx,style=lccode,#1}}
  {}

% C# comparison snippets. This book is written for engineers arriving from
% .NET, so a handful of listings show the two languages side by side.
\lstnewenvironment{csharp}[1][]
  {\lstset{language=[Sharp]C,style=lccode,#1}}
  {}

\lstnewenvironment{shellcmd}[1][]
  {\lstset{language=bash,style=lccode,numbers=none,keywordstyle=\color{lcteal}\bfseries,#1}}
  {}

\lstnewenvironment{yamlcode}[1][]
  {\lstset{style=lccode,basicstyle=\ttfamily\footnotesize,#1}}
  {}

\lstnewenvironment{jsoncode}[1][]
  {\lstset{style=lccode,basicstyle=\ttfamily\footnotesize,numbers=none,#1}}
  {}

% Terminal transcripts and tracebacks: no line numbers, no keyword colouring.
\lstnewenvironment{console}[1][]
  {\lstset{style=lccode,numbers=none,basicstyle=\ttfamily\scriptsize,
           keywordstyle=\color{black},backgroundcolor=\color{white},#1}}
  {}

% External source file — the preferred form once code lives in code/, which
% mirrors the stages of the companion mini-project repository.
% Usage: \pyfile{code/ch07/middleware.py}{Caption}{lst:label}
\newcommand{\pyfile}[3]{%
  \lstinputlisting[language=pythonx,style=lccode,caption={#2},label={#3}]{#1}%
}

% A named region of an external file, delimited in the source by
% `# --8<-- [start:name]` / `# --8<-- [end:name]` so the book and the running
% project cannot drift apart.
% Usage: \pyregion{code/ch07/middleware.py}{audit}{Caption}{lst:label}
\newcommand{\pyregion}[4]{%
  \lstinputlisting[
    language=pythonx, style=lccode,
    linerange={--8<--\ [start:#2]--- 8<--\ [end:#2]},
    includerangemarker=false,
    caption={#3}, label={#4}]{#1}%
}

% Inline literals
\newcommand{\code}[1]{\texttt{\small #1}}
\newcommand{\pkg}[1]{\texttt{\small\color{lcblue}#1}}
\newcommand{\api}[1]{\texttt{\small\color{lcteal}#1}}

% ============================================================
%  ADMONITIONS
% ============================================================

\newtcolorbox{note}[1][]{
  enhanced, breakable,
  colback=lclight, colframe=lcblue,
  boxrule=0pt, leftrule=3pt,
  arc=0pt, outer arc=0pt,
  left=8pt, right=8pt, top=6pt, bottom=6pt,
  fonttitle=\bfseries\small,
  title={Note}, #1
}

\newtcolorbox{warning}[1][]{
  enhanced, breakable,
  colback=lclight, colframe=lcred,
  boxrule=0pt, leftrule=3pt,
  arc=0pt, outer arc=0pt,
  left=8pt, right=8pt, top=6pt, bottom=6pt,
  fonttitle=\bfseries\small,
  title={Watch out}, #1
}

% For anything that is preview / alpha / likely to move. The LangChain
% ecosystem ships faster than this book does, so this box gets used a lot.
\newtcolorbox{versionbox}[1][]{
  enhanced, breakable,
  colback=lclight, colframe=lcamber,
  boxrule=0pt, leftrule=3pt,
  arc=0pt, outer arc=0pt,
  left=8pt, right=8pt, top=6pt, bottom=6pt,
  fonttitle=\bfseries\small,
  title={Version sensitive}, #1
}

% Marks a listing that has not been executed against the pinned versions.
\newtcolorbox{verifybox}[1][]{
  enhanced, breakable,
  colback=white, colframe=lcgrey,
  boxrule=0.6pt, arc=0pt, outer arc=0pt,
  left=8pt, right=8pt, top=6pt, bottom=6pt,
  fonttitle=\bfseries\small,
  title={Run this before you trust it}, #1
}

\newtcolorbox{exercisebox}[1][]{
  enhanced, breakable,
  colback=white, colframe=lcteal,
  boxrule=0pt, leftrule=3pt,
  arc=0pt, outer arc=0pt,
  left=8pt, right=8pt, top=6pt, bottom=6pt,
  fonttitle=\bfseries\small,
  title={Exercise}, #1
}

% Points at the stage of the companion repository that this section builds.
\newtcolorbox{projectbox}[1][]{
  enhanced, breakable,
  colback=white, colframe=lcamber,
  boxrule=0pt, leftrule=3pt,
  arc=0pt, outer arc=0pt,
  left=8pt, right=8pt, top=6pt, bottom=6pt,
  fonttitle=\bfseries\small,
  title={Mini-project}, #1
}

% The .NET reader is the book's assumed reader. This box carries the
% translation and nothing else, so a reader without that background can skip
% every one of them and lose no content.
\newtcolorbox{dotnetbox}[1][]{
  enhanced, breakable,
  colback=lclight, colframe=lcteal,
  boxrule=0pt, leftrule=3pt,
  arc=0pt, outer arc=0pt,
  left=8pt, right=8pt, top=6pt, bottom=6pt,
  fonttitle=\bfseries\small,
  title={If you are coming from .NET}, #1
}

% ============================================================
%  DIAGRAMS — Mermaid pipeline
%  ---------------------------------------------------------
%  \mermaidfig{key}{caption}{one-line description}
%
%  Source of truth is figures/mermaid/<key>.mmd, checked into the
%  repository. `make diagrams` renders each one to figures/diagrams/<key>.pdf
%  with mermaid-cli. If the rendered PDF is present it is included; if not, a
%  visible placeholder is drawn AND the Mermaid source is typeset verbatim, so
%  an unrendered build still shows the reader the graph.
%
%  Rendered diagrams are build output and are not committed; the .mmd source
%  is. That keeps diagrams reviewable in a pull request as text.
% ============================================================

\makeatletter
\newcommand{\listofdiagrams}{%
  \section*{Diagram manifest}%
  \addcontentsline{toc}{section}{Diagram manifest}%
  \@starttoc{dgm}%
}
\makeatother

\newcommand{\mermaidfig}[3]{%
  \addcontentsline{dgm}{subsection}{\protect\numberline{}\texttt{#1.mmd} --- #3}%
  \begin{figure}[htbp]
  \centering
  \IfFileExists{figures/diagrams/#1.pdf}{%
    \includegraphics[width=0.95\linewidth,height=0.42\textheight,keepaspectratio]{figures/diagrams/#1.pdf}%
  }{%
    \begin{tcolorbox}[
      enhanced, breakable, width=0.98\linewidth,
      colback=white, colframe=lcgrey, boxrule=0.8pt,
      arc=0pt, outer arc=0pt, left=8pt, right=8pt, top=8pt, bottom=8pt]
      {\bfseries\color{lcgrey}\small NOT RENDERED --- run \texttt{make diagrams}}\\[4pt]
      {\ttfamily\scriptsize\color{lcgrey} figures/mermaid/#1.mmd}\\[6pt]
      \IfFileExists{figures/mermaid/#1.mmd}{%
        \lstinputlisting[style=lccode,numbers=none,basicstyle=\ttfamily\scriptsize,
                         backgroundcolor=\color{lclight},frame=none,
                         keywordstyle=\color{black}]{figures/mermaid/#1.mmd}%
      }{%
        {\small\color{lcred} Mermaid source missing: \texttt{figures/mermaid/#1.mmd}}%
      }%
    \end{tcolorbox}%
  }%
  \caption{#2}\label{fig:#1}
  \end{figure}%
}

% ============================================================
%  SCREENSHOT PLACEHOLDERS
%  ---------------------------------------------------------
%  \needscreenshot{key}{caption}{what exactly to capture}
%
%  If figures/screenshots/<key>.png exists it is included. If not, a visible
%  placeholder box is drawn and the request is recorded in the manifest.
%
%  Keep capture instructions in prose. A long \code{} identifier here lands in
%  the manifest's narrow column and overflows the margin badly.
% ============================================================

\makeatletter
\newcommand{\listofshots}{%
  \section*{Screenshot manifest}%
  \addcontentsline{toc}{section}{Screenshot manifest}%
  \@starttoc{shots}%
}
\makeatother

\newcommand{\needscreenshot}[3]{%
  \addcontentsline{shots}{subsection}{\protect\numberline{}\texttt{#1.png} --- #3}%
  \begin{figure}[htbp]
  \centering
  \IfFileExists{figures/screenshots/#1.png}{%
    \includegraphics[width=0.95\linewidth]{figures/screenshots/#1.png}%
  }{%
    \begin{tcolorbox}[
      enhanced, width=0.95\linewidth,
      colback=lclight, colframe=lcamber,
      boxrule=1pt, arc=0pt, outer arc=0pt,
      borderline={0.6pt}{3pt}{lcamber, dashed},
      left=12pt, right=12pt, top=14pt, bottom=14pt]
      \centering
      {\bfseries\color{lcamber}\large SCREENSHOT NEEDED}\\[6pt]
      {\ttfamily\small figures/screenshots/#1.png}\\[10pt]
      \begin{minipage}{0.9\linewidth}
        \small\raggedright\color{lcgrey} #3
      \end{minipage}
    \end{tcolorbox}%
  }%
  \caption{#2}\label{fig:#1}
  \end{figure}%
}

% ============================================================
%  CHAPTER STUBS
%  ---------------------------------------------------------
%  \chapterstub{one-paragraph statement of what the chapter will cover}
%
%  Prints a visible marker so an unfinished chapter cannot be mistaken for a
%  finished one, and so the page count never flatters the draft.
% ============================================================
\newcommand{\chapterstub}[1]{%
  \begin{tcolorbox}[
    enhanced, breakable, width=\linewidth,
    colback=lclight, colframe=lcred,
    boxrule=1pt, arc=0pt, outer arc=0pt,
    borderline={0.6pt}{3pt}{lcred, dashed},
    left=12pt, right=12pt, top=12pt, bottom=12pt]
    {\bfseries\color{lcred}\large NOT YET WRITTEN}\\[8pt]
    \small\raggedright #1
  \end{tcolorbox}%
}

% ============================================================
%  PINNED VERSIONS — single source of truth
%  Change these here; they propagate through the whole book.
%  Verified against pypi.org on \pinnedon.
% ============================================================
\newcommand{\lcver}{1.3.14}          % langchain
\newcommand{\lccorever}{1.5.3}       % langchain-core
\newcommand{\lgver}{1.2.10}          % langgraph
\newcommand{\lgpgver}{3.1.1}         % langgraph-checkpoint-postgres
\newcommand{\lcclassicver}{1.0.8}    % langchain-classic
\newcommand{\lcanthropicver}{1.5.3}  % langchain-anthropic
\newcommand{\lsver}{0.10.15}         % langsmith
\newcommand{\mcpadapterver}{0.3.1}   % langchain-mcp-adapters
\newcommand{\deepagentsver}{0.7.3}   % deepagents
\newcommand{\pinnedon}{August 2026}  % "verified as of"
\newcommand{\pyver}{Python 3.13}     % the interpreter the book targets
\newcommand{\pymin}{Python 3.10}     % the floor LangChain 1.x requires

% The index is two-column and its entries include \texttt{} API names, which do
% not hyphenate. Ragged right keeps multi-word entries off the margin. It cannot
% help a single token wider than the column, so keep verbatim index entries under
% about 29 characters and index longer names by concept instead.
\apptocmd{\theindex}{\raggedright}{}{%
  \PackageWarning{llmbook}{Could not patch theindex for ragged-right setting}}

\makeindex
